% Exercises for Chapter 16 (numerical methods). Statements only: solutions are kept privately.
% Throughout, the Black--Scholes model is taken under the risk-neutral measure Q, N is the standard
% normal distribution function, and the finite-difference notation (a, b, c, h, k, lambda = a k / h^2,
% p_-, p_0, p_+, the theta-scheme) is that of Section 16.2 of the notes.

\begin{exercise}[When confidence intervals lie]\label{ex:num-coverage}
\easy This is a computational exercise. Take $S_0=100$, $r=5\%$, $\sigma=20\%$, $T=1$, and the digital
call that pays $\1_{\{S_T>K\}}$ at $T$, whose price is $p=e^{-rT}q$ with $q=\Q(S_T>K)=N(d_-)$
(\noteref{prop:bs-digital}). Let $Y=e^{-rT}\1_{\{S_T>K\}}$ and use the interval
\noteref{eq:num-ci} with $z=1.96$.
\begin{enumerate}[label=(\alph*)]
  \item Show that $\Var(Y)=e^{-2rT}q(1-q)$, and that the half-width of the interval divided by $p$ is
    approximately $1.96\sqrt{(1-q)/(qM)}$. How many paths make it $1\%$ for $K=100$, and for $K=200$?
  \item Take $K=200$. With probability $(1-q)^M$ all $M$ payoffs vanish. What is the interval then,
    and does it contain $p$? Deduce an upper bound for the probability that the interval contains $p$,
    for $M=10^2$, $10^3$ and $10^4$.
  \item Estimate this probability, the \emph{coverage}, for $K\in\{100,200\}$ and
    $M\in\{10^2,10^3,10^4,10^5\}$, by repeating the whole experiment many times. \emph{Hint:} the estimator
    and the sample variance depend on the payoffs only through the number of paths that end above
    $K$, whose law you know, so there is no need to simulate the paths. Compare with (b) and with
    $95\%$, and explain why this does not contradict \noteref{thm:num-mc}.
\end{enumerate}
\end{exercise}

\begin{exercise}[The explicit scheme as a tree]\label{ex:num-trinomial}
\easy Let $r\ge0$, and consider the explicit scheme for the Black--Scholes equation in log-price
\noteref{eq:num-pde}, on a grid with $h|b|\le2a$, $k(2a/h^2+c)\le1$ and $rk<1$.
\begin{enumerate}[label=(\alph*)]
  \item Write the scheme as $U^{n+1}_j=(1-rk)\big(\pi_+U^n_{j+1}+\pi_0U^n_j+\pi_-U^n_{j-1}\big)$, with the
    probabilities of \noteref{rem:num-trinomial} written out in terms of $a$, $b$, $r$, $h$ and $k$.
  \item Let $\Delta$ be a random variable equal to $h$, $0$, $-h$ with probabilities $\pi_+$, $\pi_0$,
    $\pi_-$: the move of the log-price in one period of the tree. Compute $\E[\Delta]$ and $\E[\Delta^2]$, and
    compare them with the mean and the second moment of $\log S_{t+k}-\log S_t$ in the Black--Scholes model.
    Show that both differences are of order $k^2$, whatever the choice of $h$. Compare also the discount
    factor $1-rk$ with $e^{-rk}$.
  \item Matching two moments is not enough. Compare $\E[\Delta^4]$ with the fourth moment of a normal
    variable with variance $\sigma^2k$, and explain why the scheme converges only if $h$ tends to zero
    as well as $k$. What does the tree look like when $h$ is fixed and $k\to0$?
\end{enumerate}
\end{exercise}

\begin{exercise}[A magic ratio]\label{ex:num-truncation}
\medium Consider the heat equation $\partial_\tau u=a\,\partial_{yy}u$ and the explicit scheme, with
$\lambda=ak/h^2$, and let $u$ be a solution whose partial derivatives of order at most $6$ are
continuous and bounded.
\begin{enumerate}[label=(\alph*)]
  \item Show that the truncation error is
    \[
      T^n_j=\frac{ah^2}2\Big(\lambda-\frac16\Big)\partial^4_yu(\tau_n,y_j)+O(k^2+h^4).
    \]
  \item Deduce that for $\lambda=\frac16$ the error of the scheme is $O(h^4)$, uniformly on the grid,
    when the boundary values are exact. You may use the argument of \noteref{thm:num-schemes}(iii)--(iv).
  \item This is a computational part. Take $a=\frac12$ on $[0,\pi]$ with the exact solution
    $u(\tau,y)=e^{-a\tau}\sin y$ and zero boundary values. For $J=10,20,40,80$ intervals, $h=\pi/J$,
    $k=\lambda h^2/a$, run $N=\lceil1/k\rceil$ steps and report the largest error at $\tau=Nk$, for
    $\lambda=0.4$ and $\lambda=\frac16$, with the observed orders of convergence.
  \item Which of the trinomial trees of \noteref{rem:bin-trinomial-trees} corresponds to
    $\lambda=\frac16$? (Use \noteref{rem:num-trinomial}.)
\end{enumerate}
\end{exercise}

\begin{exercise}[Crank--Nicolson and positivity]\label{ex:num-cn-positivity}
\medium Consider the heat equation $\partial_\tau u=a\,\partial_{yy}u$ and the Crank--Nicolson scheme
with $\lambda=ak/h^2$, written as $BU^{n+1}=CU^n$ with $B=I-\frac k2A$ and $C=I+\frac k2A$.
\begin{enumerate}[label=(\alph*)]
  \item On an interval with zero boundary values, show that if $\lambda\le1$, then $U^n\ge0$ implies
    $U^{n+1}\ge0$. \emph{Hint:} look at the weights of $C$, and show that $Bz\ge0$ implies $z\ge0$, as
    in the proof of \noteref{prop:num-american-grid}(iii).
  \item Now work on the infinite grid $h\Z$, with bounded sequences. Show that $C=2I-B$, so that one
    step maps $U$ to $2B^{-1}U-U$. Find $\rho\in(0,1)$ and $C_\lambda>0$ such that $z_j=C_\lambda\rho^{|j|}$
    solves $Bz=e$, where $e_0=1$ and $e_j=0$ for $j\ne0$ (you may admit that it is the only bounded
    solution). Deduce that the spike $e$ is mapped to a sequence whose value at $0$ is
    $2/\sqrt{1+2\lambda}-1$.
  \item Conclude that, on the infinite grid, the step maps every nonnegative bounded sequence to a
    nonnegative one if and only if $\lambda\le\frac32$.
  \item The second differences of the call payoff on a grid whose nodes include $\log K$ vanish
    except at $\log K$: they form a spike. Since the scheme commutes with differences, use (b) to
    explain the sign of the Crank--Nicolson gamma at the strike after one step, for
    $\lambda\approx22$ as in \noteref{fig:num-convergence}, and the oscillations that follow.
\end{enumerate}
\end{exercise}

\begin{exercise}[Richardson extrapolation of a biased estimator]\label{ex:num-richardson}
\medium Work in the setting of \noteref{exa:num-power}: $Y_h=e^{-rT}(S^h_T)^2$, with the Euler scheme on
$S$ and $n=T/h$ steps, and $p=S_0^2e^{(r+\sigma^2)T}$.
\begin{enumerate}[label=(\alph*)]
  \item From the closed form of $\E[Y_h]$, show that the bias has an expansion
    $\beta(h)=\beta_1h+\beta_2h^2+O(h^3)$, and compute $\beta_1$ and $\beta_2$.
  \item Let $Z=2\hat p_{M_1,h/2}-\hat p_{M_2,h}$, where the two Monte Carlo estimators, with steps $h/2$ and
    $h$, use independent samples. Show that $\E[Z]-p=-\frac12\beta_2h^2+O(h^3)$, and compute $\Var(Z)$.
  \item Show that, for small $\varepsilon$, a root mean square error $\varepsilon$ can be reached with $Z$
    at a cost of order $\varepsilon^{-5/2}$, against $\varepsilon^{-3}$ in \noteref{prop:num-mc-cost}.
  \item Take $S_0=1$, $r=5\%$, $\sigma=40\%$, $T=1$ and $\varepsilon=10^{-3}$. Using the exact bias and
    variance (the functions of \coderef{code/src/ieor222/numerics.py} compute them), find the cheapest
    choice of $n$ and $M$ for the plain estimator, and of $n$, $M_1$ and $M_2$ for $Z$, and compare the
    costs, counted in normal draws. Compare also with the choice of \noteref{prop:num-mc-cost}(i) with
    $c_2=0.039$ and $v_2=1.37$.
\end{enumerate}
\end{exercise}

\begin{exercise}[An up-and-out call on a grid]\label{ex:num-barrier}
\medium This is a computational exercise. Take $S_0=K=100$, a barrier $B=120$, $r=5\%$, $\sigma=20\%$ and
$T=1$; the up-and-out call costs $1.1761$ (\noteref{prop:exo-barrier}). As in
\noteref{rem:exo-barrier-reading}, accept that while the option is alive its price is $v(t,S_t)$,
where $v$ solves the Black--Scholes equation on $[0,T)\times(0,B)$ with $v(t,B)=0$ and
$v(T,x)=(x-K)^+$ for $x<B$.
\begin{enumerate}[label=(\alph*)]
  \item Write the problem in log-price on a grid whose top node is $\log B$ and on which $\log S_0$ is a
    node: $h=\log(B/S_0)/m_0$ for an integer $m_0$, extending at least $6\sigma\sqrt T$ below $\log S_0$.
    What should the lower boundary value be?
  \item Solve it by Crank--Nicolson with $k\approx h/\sigma$, with and without Rannacher's start, for
    $m_0=5,10,\dots,160$, and report the errors at $S_0$ and the observed orders. Where is the payoff
    not smooth, and why does Rannacher's start matter even though the payoff is continuous at the
    strike?
  \item Now let the barrier be monitored only at the dates $iT/m$, $1\le i\le m$. Solve the equation
    between two dates on a grid that extends above $B$ (with the boundary values of a call at the top),
    and at each monitoring date set the solution to $0$ above $B$. Compute the price for $m=4,12,52,250$.
    Compare the convergence in $m_0$ when the value at the node $\log B$ itself is set to $0$, and when
    it is halved (the average of the solution over the cell around that node). Compare the prices with
    the continuously monitored price with the barrier shifted to $Be^{\beta\sigma\sqrt{T/m}}$,
    $\beta\approx0.5826$ (\noteref{rem:exo-discrete}).
\end{enumerate}
\end{exercise}

\begin{exercise}[The exercise region of the American put]\label{ex:num-american-region}
\hard Let $W$ be a Brownian motion with its natural filtration, $K>0$, $\sigma>0$, $r\ge0$, and
$R_s=\exp\big(\sigma W_s+(r-\frac{\sigma^2}2)s\big)$. For $\theta>0$, let $\mathcal T_\theta$ be the set of
stopping times with values in $[0,\theta]$, and
\[
  V(\theta,x)=\sup_{\varrho\in\mathcal T_\theta}\E\big[e^{-r\varrho}\big(K-xR_\varrho\big)^+\big],\qquad x>0,
\]
so that the price \noteref{eq:num-american} is $v(t,x)=V(T-t,x)$, by the restart property of Brownian
motion. You may use optional sampling for continuous martingales and bounded stopping times: if
$\varrho\le\varrho'$ are in $\mathcal T_\theta$, then $\E[e^{-r\varrho'}R_{\varrho'}\mid\F_\varrho]=e^{-r\varrho}R_\varrho$.
Let $E_\theta=\{x>0:V(\theta,x)=(K-x)^+\}$ be the exercise region with time $\theta$ to maturity.
\begin{enumerate}[label=(\alph*)]
  \item Show that $x\mapsto V(\theta,x)$ is convex and nonincreasing, and that
    $V(\theta,y)\le V(\theta,x)+x-y$ for $0<y<x$.
  \item Show that if $x\in E_\theta$ and $0<y<x$, then $y\in E_\theta$; that $E_\theta$ is closed in
    $(0,\infty)$; and that $E_\theta\cap[K,\infty)=\emptyset$. Deduce that $E_\theta$ is empty or of the form
    $(0,b]$ with $b<K$.
  \item Let $r>0$ and $L^*=2rK/(2r+\sigma^2)$. Using \noteref{prop:exo-perpetual-put}, show that
    $(0,L^*]\subseteq E_\theta$, so that $E_\theta=(0,b(\theta)]$ with $L^*\le b(\theta)<K$.
  \item Show that $\theta\mapsto V(\theta,x)$ is nondecreasing, and deduce that $b$ is nonincreasing in
    $\theta$: in calendar time $t=T-\theta$, the exercise boundary rises towards maturity.
  \item Let $r=0$. Show that $V(\theta,x)=\E[(K-xR_\theta)^+]>(K-x)^+$ for every $x>0$, so that
    $E_\theta=\emptyset$: the American put is never exercised early.
\end{enumerate}
\end{exercise}
